% Options for packages loaded elsewhere
% Options for packages loaded elsewhere
\PassOptionsToPackage{unicode}{hyperref}
\PassOptionsToPackage{hyphens}{url}
%
\documentclass[
  ignorenonframetext,
  aspectratio=169,
]{beamer}
\newif\ifbibliography
\usepackage{pgfpages}
\setbeamertemplate{caption}[numbered]
\setbeamertemplate{caption label separator}{: }
\setbeamercolor{caption name}{fg=normal text.fg}
\beamertemplatenavigationsymbolsempty
% remove section numbering
\setbeamertemplate{part page}{
  \centering
  \begin{beamercolorbox}[sep=16pt,center]{part title}
    \usebeamerfont{part title}\insertpart\par
  \end{beamercolorbox}
}
\setbeamertemplate{section page}{
  \centering
  \begin{beamercolorbox}[sep=12pt,center]{section title}
    \usebeamerfont{section title}\insertsection\par
  \end{beamercolorbox}
}
\setbeamertemplate{subsection page}{
  \centering
  \begin{beamercolorbox}[sep=8pt,center]{subsection title}
    \usebeamerfont{subsection title}\insertsubsection\par
  \end{beamercolorbox}
}
% Prevent slide breaks in the middle of a paragraph
\widowpenalties 1 10000
\raggedbottom
\AtBeginPart{
  \frame{\partpage}
}
\AtBeginSection{
  \ifbibliography
  \else
    \frame{\sectionpage}
  \fi
}
\AtBeginSubsection{
  \frame{\subsectionpage}
}
\usepackage{iftex}
\ifPDFTeX
  \usepackage[T1]{fontenc}
  \usepackage[utf8]{inputenc}
  \usepackage{textcomp} % provide euro and other symbols
\else % if luatex or xetex
  \usepackage{unicode-math} % this also loads fontspec
  \defaultfontfeatures{Scale=MatchLowercase}
  \defaultfontfeatures[\rmfamily]{Ligatures=TeX,Scale=1}
\fi
\usepackage{lmodern}

\usetheme[]{Boadilla}
\usecolortheme[]{default}
\usefonttheme[]{professionalfonts}
\ifPDFTeX\else
  % xetex/luatex font selection
\fi
% Use upquote if available, for straight quotes in verbatim environments
\IfFileExists{upquote.sty}{\usepackage{upquote}}{}
\IfFileExists{microtype.sty}{% use microtype if available
  \usepackage[]{microtype}
  \UseMicrotypeSet[protrusion]{basicmath} % disable protrusion for tt fonts
}{}
\makeatletter
\@ifundefined{KOMAClassName}{% if non-KOMA class
  \IfFileExists{parskip.sty}{%
    \usepackage{parskip}
  }{% else
    \setlength{\parindent}{0pt}
    \setlength{\parskip}{6pt plus 2pt minus 1pt}}
}{% if KOMA class
  \KOMAoptions{parskip=half}}
\makeatother


\usepackage{longtable,booktabs,array}
\newcounter{none} % for unnumbered tables
\usepackage{calc} % for calculating minipage widths
\usepackage{caption}
% Make caption package work with longtable
\makeatletter
\def\fnum@table{\tablename~\thetable}
\makeatother
\usepackage{graphicx}
\makeatletter
\newsavebox\pandoc@box
\newcommand*\pandocbounded[1]{% scales image to fit in text height/width
  \sbox\pandoc@box{#1}%
  \Gscale@div\@tempa{\textheight}{\dimexpr\ht\pandoc@box+\dp\pandoc@box\relax}%
  \Gscale@div\@tempb{\linewidth}{\wd\pandoc@box}%
  \ifdim\@tempb\p@<\@tempa\p@\let\@tempa\@tempb\fi% select the smaller of both
  \ifdim\@tempa\p@<\p@\scalebox{\@tempa}{\usebox\pandoc@box}%
  \else\usebox{\pandoc@box}%
  \fi%
}
% Set default figure placement to htbp
\def\fps@figure{htbp}
\makeatother





\setlength{\emergencystretch}{3em} % prevent overfull lines

\providecommand{\tightlist}{%
  \setlength{\itemsep}{0pt}\setlength{\parskip}{0pt}}



 


\makeatletter
\@ifpackageloaded{caption}{}{\usepackage{caption}}
\AtBeginDocument{%
\ifdefined\contentsname
  \renewcommand*\contentsname{Table of contents}
\else
  \newcommand\contentsname{Table of contents}
\fi
\ifdefined\listfigurename
  \renewcommand*\listfigurename{List of Figures}
\else
  \newcommand\listfigurename{List of Figures}
\fi
\ifdefined\listtablename
  \renewcommand*\listtablename{List of Tables}
\else
  \newcommand\listtablename{List of Tables}
\fi
\ifdefined\figurename
  \renewcommand*\figurename{Figure}
\else
  \newcommand\figurename{Figure}
\fi
\ifdefined\tablename
  \renewcommand*\tablename{Table}
\else
  \newcommand\tablename{Table}
\fi
}
\@ifpackageloaded{float}{}{\usepackage{float}}
\floatstyle{ruled}
\@ifundefined{c@chapter}{\newfloat{codelisting}{h}{lop}}{\newfloat{codelisting}{h}{lop}[chapter]}
\floatname{codelisting}{Listing}
\newcommand*\listoflistings{\listof{codelisting}{List of Listings}}
\makeatother
\makeatletter
\makeatother
\makeatletter
\@ifpackageloaded{caption}{}{\usepackage{caption}}
\@ifpackageloaded{subcaption}{}{\usepackage{subcaption}}
\makeatother

\usepackage{bookmark}
\IfFileExists{xurl.sty}{\usepackage{xurl}}{} % add URL line breaks if available
\urlstyle{same}
\hypersetup{
  pdftitle={Lecture 1: Introduction to Stochastic Processes},
  hidelinks,
  pdfcreator={LaTeX via pandoc}}


\title{\texorpdfstring{Lecture 1: Introduction to Stochastic
Processes}{Lecture 1: Introduction to Stochastic Processes}}
\author{}
\date{}

\begin{document}
\frame{\titlepage}


\begin{frame}{Introduction to Stochastic Processes}
\protect\phantomsection\label{introduction-to-stochastic-processes}
\vspace{0.5cm}

\textbf{Course:} Stochastic Processes

\textbf{Reference:} IFoA CS2

\vfill

Department of Mathematics /Actuarial Science Program
\end{frame}

\begin{frame}{Learning Outcomes}
\protect\phantomsection\label{learning-outcomes}
After this lecture, you should be able to

\begin{enumerate}
\item
  \textbf{Explain} why stochastic processes are used to model systems
  that evolve over time under uncertainty.
\item
  \textbf{Distinguish} between a random variable and a stochastic
  process.
\item
  \textbf{Identify} the stochastic process, time domain, and state space
  for simple actuarial applications.
\item
  \textbf{Define} a stochastic process using the notation \[
  \{X_t:t\in T\}.
  \]
\item
  \textbf{Classify} a stochastic process according to its time domain
  and state space.
\end{enumerate}
\end{frame}

\begin{frame}{Why Study Stochastic Processes?}
\protect\phantomsection\label{why-study-stochastic-processes}
Many quantities in actuarial science change over time.

Unlike deterministic models, their future values are \textbf{uncertain}.

Some common examples are

\begin{itemize}
\tightlist
\item
  \textbf{Insurance premiums} --- change according to a policyholder's
  claim history.
\item
  \textbf{Claim arrivals} --- occur at random times.
\item
  \textbf{Survival of a life} --- depends on an uncertain future
  lifetime.
\item
  \textbf{Stock prices} --- fluctuate continuously over time.
\end{itemize}

\begin{quote}
\textbf{Observation:} We need mathematical models that describe both
\textbf{time evolution} and \textbf{uncertainty}.
\end{quote}
\end{frame}

\begin{frame}{Examples from Actuarial Science}
\protect\phantomsection\label{examples-from-actuarial-science}
{\def\LTcaptype{none} % do not increment counter
\begin{longtable}[]{@{}ll@{}}
\toprule\noalign{}
Application & Quantity of Interest \\
\midrule\noalign{}
\endhead
No-Claim Discount (NCD) & Discount level of a policyholder \\
General insurance & Number of claims received \\
Life insurance & Survival status of an individual \\
Finance & Price of a financial asset \\
\bottomrule\noalign{}
\end{longtable}
}

Each of these quantities

\begin{itemize}
\tightlist
\item
  \textbf{Changes over time} rather than remaining fixed.
\item
  \textbf{Cannot be predicted with certainty} because future outcomes
  are random.
\item
  \textbf{Can be modelled} using a stochastic process.
\end{itemize}
\end{frame}

\begin{frame}{Common Characteristics}
\protect\phantomsection\label{common-characteristics}
Although the applications are different, they share the same modelling
framework.

\begin{itemize}
\tightlist
\item
  \textbf{State variable} --- describes the current status of the
  system.
\item
  \textbf{Time index} --- specifies when the system is observed.
\item
  \textbf{Uncertainty} --- future states cannot be known exactly.
\item
  \textbf{Evolution} --- the state changes as time progresses.
\end{itemize}

These common features motivate the formal study of \textbf{stochastic
processes}.

\begin{quote}
\textbf{Key Idea:} A stochastic process is a mathematical model for a
random quantity that evolves over time.
\end{quote}
\end{frame}

\begin{frame}{Random Variables and Stochastic Processes}
\protect\phantomsection\label{random-variables-and-stochastic-processes}
Many random quantities are observed repeatedly over time.

Examples include

\begin{itemize}
\tightlist
\item
  Daily stock prices.
\item
  Annual No-Claim Discount (NCD) levels.
\item
  Number of insurance claims received.
\item
  Survival status of an individual.
\end{itemize}

Instead of studying \textbf{one} random variable, we study a
\textbf{collection of random variables}, one for each time point.
\end{frame}

\begin{frame}{Random Variable}
\protect\phantomsection\label{random-variable}
A \textbf{random variable} describes the outcome of a random experiment
at a \textbf{single time point}.

Examples

\begin{itemize}
\tightlist
\item
  \(X\) = claim amount from a policy.
\item
  \(Y\) = lifetime of an individual.
\item
  \(Z\) = number of claims reported today.
\end{itemize}

A random variable provides information about \textbf{one observation}
only.

\begin{quote}
\textbf{Observation:} A single random variable does not describe how a
quantity changes over time.
\end{quote}
\end{frame}

\begin{frame}{From a Random Variable to a Stochastic Process}
\protect\phantomsection\label{from-a-random-variable-to-a-stochastic-process}
Suppose the NCD level of a policyholder is recorded every year.

Instead of defining one random variable,

\[
X,
\]

we define a sequence of random variables,

\[
X_0,\;X_1,\;X_2,\;\ldots
\]

where

\begin{itemize}
\tightlist
\item
  \(X_0\) is the initial NCD level.
\item
  \(X_1\) is the NCD level after one year.
\item
  \(X_2\) is the NCD level after two years, \(\ldots\)
\end{itemize}

Each random variable describes the system at a different point in time.
\end{frame}

\begin{frame}{A Collection of Random Variables}
\protect\phantomsection\label{a-collection-of-random-variables}
The sequence

\[
X_0,\;X_1,\;X_2,\;\ldots
\]

can be written more compactly as

\[
\{X_t:t\in T\},
\]

where

\begin{itemize}
\tightlist
\item
  \(X_t\) is the random variable observed at time \(t\).
\item
  \(T\) is the collection of time points.
\end{itemize}

This notation represents a \textbf{stochastic process}.
\end{frame}

\begin{frame}{Understanding the Notation}
\protect\phantomsection\label{understanding-the-notation}
The notation

\[
\{X_t:t\in T\}
\]

has two important components.

\begin{itemize}
\tightlist
\item
  \textbf{\(X_t\)} --- the random variable observed at time \(t\).
\item
  \textbf{\(t\)} --- the time index.
\item
  \textbf{\(T\)} --- the set of all time points.
\end{itemize}

The value of \(T\) depends on the application.

For example,

\begin{itemize}
\tightlist
\item
  \(T=\{0,1,2,\ldots\}\) for annual observations.
\item
  \(T=[0,\infty)\) for continuous-time observations.
\end{itemize}
\end{frame}

\begin{frame}{Random Variable vs Stochastic Process}
\protect\phantomsection\label{random-variable-vs-stochastic-process}
{\def\LTcaptype{none} % do not increment counter
\begin{longtable}[]{@{}
  >{\raggedright\arraybackslash}p{(\linewidth - 2\tabcolsep) * \real{0.4595}}
  >{\raggedright\arraybackslash}p{(\linewidth - 2\tabcolsep) * \real{0.5405}}@{}}
\toprule\noalign{}
\begin{minipage}[b]{\linewidth}\raggedright
Random Variable
\end{minipage} & \begin{minipage}[b]{\linewidth}\raggedright
Stochastic Process
\end{minipage} \\
\midrule\noalign{}
\endhead
Describes one random outcome. & Describes the evolution of a random
quantity over time. \\
Represented by a single variable such as \(X\). & Represented by a
collection \(\{X_t:t\in T\}\). \\
No time index is required. & Every random variable is indexed by
time. \\
\bottomrule\noalign{}
\end{longtable}
}

\begin{quote}
\textbf{Key Idea:} A stochastic process is \textbf{not} a single random
variable. It is a collection of random variables indexed by time.
\end{quote}
\end{frame}

\begin{frame}{Looking Ahead}
\protect\phantomsection\label{looking-ahead}
In the next section, we will study several stochastic processes from
actuarial science.

For each example, we will identify

\begin{itemize}
\tightlist
\item
  the \textbf{stochastic process},
\item
  the \textbf{time domain},
\item
  the \textbf{state space}, and
\item
  the \textbf{type of stochastic process}.
\end{itemize}

These examples provide the foundation for the formal definition
introduced later.
\end{frame}

\begin{frame}{Example 1: No-Claim Discount (NCD) System}
\protect\phantomsection\label{example-1-no-claim-discount-ncd-system}
A \textbf{No-Claim Discount (NCD)} system rewards policyholders who do
not make claims by offering lower insurance premiums.

The NCD level is reviewed at the end of each policy year.

Depending on the policyholder's claim history, the NCD level may

\begin{itemize}
\tightlist
\item
  \textbf{Increase} after a claim-free year.
\item
  \textbf{Decrease} after one or more claims.
\item
  \textbf{Remain unchanged} under the insurer's rules.
\end{itemize}

\begin{quote}
\textbf{Observation:} The NCD level changes over time and its future
value is uncertain.
\end{quote}
\end{frame}

\begin{frame}{Modelling the NCD System}
\protect\phantomsection\label{modelling-the-ncd-system}
Suppose the insurer records the NCD level once every policy year.

Let

\[
X_n=\text{NCD level after the }n\text{th policy year},
\]

where

\begin{itemize}
\tightlist
\item
  \(n=0,1,2,\ldots\) denotes the policy year.
\item
  \(X_n\) is a random variable representing the NCD level at year \(n\).
\end{itemize}

The collection

\[
\{X_n:n=0,1,2,\ldots\}
\]

forms a \textbf{stochastic process}.
\end{frame}

\begin{frame}{Why is \(X_n\) Random?}
\protect\phantomsection\label{why-is-x_n-random}
At the beginning of the year, the insurer does \textbf{not} know whether
the policyholder will make a claim.

Consequently,

\begin{itemize}
\tightlist
\item
  a claim-free year may increase the NCD level,
\item
  one or more claims may reduce the NCD level, and
\item
  the future NCD level cannot be predicted with certainty.
\end{itemize}

Therefore, each \(X_n\) is a \textbf{random variable}.

\begin{quote}
\textbf{Key Idea:} The randomness comes from the policyholder's future
claim experience.
\end{quote}
\end{frame}

\begin{frame}{Time Domain and State Space}
\protect\phantomsection\label{time-domain-and-state-space}
The NCD process is

\[
\{X_n:n=0,1,2,\ldots\}.
\]

Its two main components are

\begin{itemize}
\item
  \textbf{Time domain}

  \[
  T=\{0,1,2,\ldots\},
  \]

  since the NCD level is observed once every policy year.
\item
  \textbf{State space}

  The set of all possible NCD levels defined by the insurer.

  For example,

  \[
  S=\{0,1,2,3,4\}.
  \]
\end{itemize}

The exact state space depends on the insurer's NCD scheme.
\end{frame}

\begin{frame}{Classification of the NCD Process}
\protect\phantomsection\label{classification-of-the-ncd-process}
The NCD process has the following characteristics.

\begin{itemize}
\item
  \textbf{Discrete-time}

  Observations are made at distinct time points (policy years).
\item
  \textbf{Discrete-state}

  Only a finite number of NCD levels are possible.
\end{itemize}

Hence,

\[
\boxed{\text{The NCD system is a discrete-time, discrete-state stochastic process.}}
\]

This type of process is studied in detail using \textbf{discrete-time
Markov chains} later in the course.
\end{frame}

\begin{frame}{Example 2: Insurance Claim Arrivals}
\protect\phantomsection\label{example-2-insurance-claim-arrivals}
Insurance companies receive claims throughout the day.

Examples include

\begin{itemize}
\tightlist
\item
  Motor insurance claims.
\item
  Property insurance claims.
\item
  Health insurance claims.
\end{itemize}

Unlike the NCD example, claims may occur \textbf{at any instant}.

Typical actuarial questions include

\begin{itemize}
\tightlist
\item
  How many claims will be reported today?
\item
  What is the probability of receiving more than 100 claims this week?
\item
  How long until the next claim arrives?
\end{itemize}

\begin{quote}
\textbf{Observation:} We require a model that describes random events
occurring continuously over time.
\end{quote}
\end{frame}

\begin{frame}{Modelling Claim Arrivals}
\protect\phantomsection\label{modelling-claim-arrivals}
Let

\[
N_t=\text{number of claims received by time }t,
\]

where

\begin{itemize}
\tightlist
\item
  \(t\ge0\) denotes time.
\item
  \(N_t\) counts all claims reported from time \(0\) up to time \(t\).
\end{itemize}

The collection

\[
\{N_t:t\ge0\}
\]

forms a \textbf{stochastic process}.
\end{frame}

\begin{frame}{Why is \(N_t\) Random?}
\protect\phantomsection\label{why-is-n_t-random}
At the present time, the insurer cannot predict

\begin{itemize}
\tightlist
\item
  when the next claim will occur,
\item
  how many claims will be received during the next hour or day, or
\item
  the total number of claims by a future time.
\end{itemize}

Consequently,

\begin{itemize}
\tightlist
\item
  the value of \(N_t\) is uncertain for every future time \(t\), and
\item
  each \(N_t\) is a \textbf{random variable}.
\end{itemize}

\begin{quote}
\textbf{Key Idea:} The randomness comes from the unpredictable arrival
times of future claims.
\end{quote}
\end{frame}

\begin{frame}{Time Domain and State Space}
\protect\phantomsection\label{time-domain-and-state-space-1}
The claim-count process is

\[
\{N_t:t\ge0\}.
\]

Its two main components are

\begin{itemize}
\item
  \textbf{Time domain} (claims occur at any instant)

  \[
  T=[0,\infty),
  \]
\item
  \textbf{State space} (count of non-negative claims)

  \[
  S=\{0,1,2,\ldots\},
  \]
\end{itemize}
\end{frame}

\begin{frame}{Classification of the Claim-Count Process}
\protect\phantomsection\label{classification-of-the-claim-count-process}
The process has the following characteristics.

\begin{itemize}
\item
  \textbf{Continuous-time}

  Time is measured continuously.
\item
  \textbf{Discrete-state}

  The process records the number of claims, which is always a whole
  number.
\end{itemize}

Hence,

\[
\boxed{\text{The claim-count process is a continuous-time, discrete-state stochastic process.}}
\]

A fundamental example of this type of process is the \textbf{Poisson
process}, which will be studied in detail later in the course.
\end{frame}

\begin{frame}{Example 3: Survival of a Life}
\protect\phantomsection\label{example-3-survival-of-a-life}
Life insurance contracts depend on the future lifetime of the insured
individual.

Actuaries are interested in questions such as

\begin{itemize}
\tightlist
\item
  What is the probability that a life aged 70 survives for another year?
\item
  What is the probability that death occurs before age 71?
\item
  How does the survival status change over time?
\end{itemize}

\begin{quote}
\textbf{Observation:} The survival status changes over time, but the
future lifetime is uncertain.
\end{quote}
\end{frame}

\begin{frame}{Modelling Survival Status}
\protect\phantomsection\label{modelling-survival-status}
Suppose we observe an individual continuously over time.

Define

\[
X_t=
\begin{cases}
1, & \text{if the individual is alive at time }t,\\
0, & \text{if the individual has died by time }t.
\end{cases}
\]

where

\begin{itemize}
\tightlist
\item
  \(t\ge0\) denotes the time since observation begins.
\item
  \(X_t\) represents the survival status at time \(t\).
\end{itemize}

The collection

\[
\{X_t:t\ge0\}
\]

forms a \textbf{stochastic process}.
\end{frame}

\begin{frame}{Why is \(X_t\) Random?}
\protect\phantomsection\label{why-is-x_t-random}
At the present time, we do not know

\begin{itemize}
\tightlist
\item
  when death will occur,
\item
  whether the individual will still be alive at a future time, or
\item
  how long the individual will survive.
\end{itemize}

Therefore,

\begin{itemize}
\tightlist
\item
  the value of \(X_t\) is uncertain for every future time \(t\), and
\item
  each \(X_t\) is a \textbf{random variable}.
\end{itemize}

\begin{quote}
\textbf{Key Idea:} The randomness arises from the uncertain future
lifetime of the individual.
\end{quote}
\end{frame}

\begin{frame}{Time Domain and State Space}
\protect\phantomsection\label{time-domain-and-state-space-2}
The survival process is

\[
\{X_t:t\ge0\}.
\]

Its two main components are

\begin{itemize}
\item
  \textbf{Time domain} (survival can be monitored continuously over
  time)

  \[
  T=[0,\infty),
  \]
\item
  \textbf{State space}

  \[
  S=\{0,1\},
  \]

  where \(1\) represents \textbf{alive}, and \(0\) represents
  \textbf{dead}.
\end{itemize}
\end{frame}

\begin{frame}{Classification of the Survival Process}
\protect\phantomsection\label{classification-of-the-survival-process}
The process has the following characteristics.

\begin{itemize}
\item
  \textbf{Continuous-time}

  The survival status is defined at every instant.
\item
  \textbf{Discrete-state}

  Only two states are possible: alive or dead.
\end{itemize}

Hence,

\[
\boxed{\text{The survival process is a continuous-time, discrete-state stochastic process.}}
\]

This type of process is commonly modelled using \textbf{continuous-time
Markov processes}, which will be studied later in the course.
\end{frame}

\begin{frame}{From Examples to a General Model}
\protect\phantomsection\label{from-examples-to-a-general-model}
Although the applications are different, they have the same mathematical
structure.

{\def\LTcaptype{none} % do not increment counter
\begin{longtable}[]{@{}ll@{}}
\toprule\noalign{}
Application & Stochastic Process \\
\midrule\noalign{}
\endhead
NCD system & \(\{X_n:n=0,1,2,\ldots\}\) \\
Claim arrivals & \(\{N_t:t\ge0\}\) \\
Survival status & \(\{X_t:t\ge0\}\) \\
\bottomrule\noalign{}
\end{longtable}
}

Each example consists of

\begin{itemize}
\tightlist
\item
  a \textbf{collection of random variables},
\item
  indexed by \textbf{time}, and
\item
  describing the \textbf{evolution of a system}.
\end{itemize}

\begin{quote}
\textbf{Observation:} We now introduce a general definition that applies
to all stochastic processes.
\end{quote}
\end{frame}

\begin{frame}{Formal Definition}
\protect\phantomsection\label{formal-definition}
A \textbf{stochastic process} is a collection of random variables

\[
\{X_t:t\in T\},
\]

defined on a common probability space

\[
(\Omega,\mathcal{F},P),
\]

where

\begin{itemize}
\tightlist
\item
  \textbf{\(X_t\)} is the random variable observed at time \(t\).
\item
  \textbf{\(T\)} is the index set (time domain).
\end{itemize}

Each random variable describes the state of the system at a particular
time.
\end{frame}

\begin{frame}{Understanding the Definition}
\protect\phantomsection\label{understanding-the-definition}
The notation

\[
\{X_t:t\in T\}
\]

contains three important components.

\begin{itemize}
\item
  \textbf{Random variables} \(X_t\)

  Describe the state of the system at time \(t\).
\item
  \textbf{Time index} \(t\)

  Identifies a particular observation time.
\item
  \textbf{Index set} \(T\)

  Contains all time points for which the process is defined.
\end{itemize}

Different applications may use different index sets.
\end{frame}

\begin{frame}{Index Set vs State Space}
\protect\phantomsection\label{index-set-vs-state-space}
Two concepts are easily confused.

{\def\LTcaptype{none} % do not increment counter
\begin{longtable}[]{@{}
  >{\raggedright\arraybackslash}p{(\linewidth - 2\tabcolsep) * \real{0.4800}}
  >{\raggedright\arraybackslash}p{(\linewidth - 2\tabcolsep) * \real{0.5200}}@{}}
\toprule\noalign{}
\begin{minipage}[b]{\linewidth}\raggedright
Index Set
\end{minipage} & \begin{minipage}[b]{\linewidth}\raggedright
State Space
\end{minipage} \\
\midrule\noalign{}
\endhead
Specifies \textbf{when} the process is observed. & Specifies
\textbf{what values} the process can take. \\
Usually denoted by \(T\). & Usually denoted by \(S\). \\
\bottomrule\noalign{}
\end{longtable}
}

Examples

{\def\LTcaptype{none} % do not increment counter
\begin{longtable}[]{@{}lll@{}}
\toprule\noalign{}
Process & Index Set \(T\) & State Space \(S\) \\
\midrule\noalign{}
\endhead
NCD & \(\{0,1,2,\ldots\}\) & \(\{0,1,2,3,4\}\) \\
Claim arrivals & \([0,\infty)\) & \(\{0,1,2,\ldots\}\) \\
Survival & \([0,\infty)\) & \(\{0,1\}\) \\
\bottomrule\noalign{}
\end{longtable}
}

\begin{quote}
\textbf{Key Idea:} The index set describes \textbf{time}, whereas the
state space describes \textbf{possible values}.
\end{quote}
\end{frame}

\begin{frame}{Looking Ahead}
\protect\phantomsection\label{looking-ahead-1}
The formal definition introduces several important components.

\begin{itemize}
\tightlist
\item
  \textbf{Probability space} \((\Omega,\mathcal{F},P)\) --- defines the
  underlying random experiment.
\item
  \textbf{Index set} \(T\) --- specifies when observations are made.
\item
  \textbf{State space} \(S\) --- contains all possible values of the
  process.
\item
  \textbf{Sample path} --- represents one possible evolution of the
  process.
\end{itemize}

In the next section, we examine each of these components in more detail.
\end{frame}

\begin{frame}{Components of a Stochastic Process}
\protect\phantomsection\label{components-of-a-stochastic-process}
A stochastic process

\[
\{X_t:t\in T\}
\]

contains several important components.

Throughout this course, we will identify

\begin{itemize}
\tightlist
\item
  the \textbf{probability space} \((\Omega,\mathcal{F},P)\),
\item
  the \textbf{index set} \(T\),
\item
  the \textbf{state space} \(S\), and
\item
  the \textbf{sample path}.
\end{itemize}

Understanding these components helps us describe and analyse any
stochastic process.
\end{frame}

\begin{frame}{Probability Space}
\protect\phantomsection\label{probability-space}
Every stochastic process is defined on a probability space

\[
(\Omega,\mathcal{F},P),
\]

where

\begin{itemize}
\tightlist
\item
  \textbf{Sample space} \(\Omega\) --- contains all possible outcomes of
  the experiment.
\item
  \textbf{Sigma-algebra} \(\mathcal{F}\) --- specifies the events of
  interest.
\item
  \textbf{Probability measure} \(P\) --- assigns probabilities to those
  events.
\end{itemize}

The probability space provides the mathematical foundation for the
stochastic process.

\begin{quote}
\textbf{Remark:} In this course, we will usually focus on the stochastic
process itself rather than the details of the probability space.
\end{quote}
\end{frame}

\begin{frame}{Index Set (Time Domain)}
\protect\phantomsection\label{index-set-time-domain}
The \textbf{index set} specifies the time points at which the process is
defined.

Common choices are

\begin{itemize}
\item
  \textbf{Discrete time}: observations are made at distinct time points

  \[
  T=\{0,1,2,\ldots\},
  \]
\item
  \textbf{Continuous time}: the process is defined for every instant.

  \[
  T=[0,\infty),
  \]
\end{itemize}

Examples (1) NCD system: yearly observations, (2) Claim arrivals:
continuous time, and (3) Survival process: continuous time.
\end{frame}

\begin{frame}{State Space}
\protect\phantomsection\label{state-space}
The \textbf{state space} is the set of all possible values that the
process can take.

Examples

{\def\LTcaptype{none} % do not increment counter
\begin{longtable}[]{@{}ll@{}}
\toprule\noalign{}
Process & State Space \\
\midrule\noalign{}
\endhead
NCD system & \(\{0,1,2,3,4\}\) \\
Claim arrivals & \(\{0,1,2,\ldots\}\) \\
Survival process & \(\{0,1\}\) \\
\bottomrule\noalign{}
\end{longtable}
}

The choice of state space depends on the quantity being modelled.

\begin{quote}
\textbf{Key Idea:} The state space describes \textbf{what values} the
process can assume, not \textbf{when} they occur.
\end{quote}
\end{frame}

\begin{frame}{Sample Paths}
\protect\phantomsection\label{sample-paths}
A \textbf{sample path} is one possible realisation of a stochastic
process.

It records how the state of the process changes over time.

Examples

\begin{itemize}
\tightlist
\item
  An NCD level that changes over successive policy years.
\item
  A cumulative claim count observed over time.
\item
  The survival status of an individual over time.
\end{itemize}

Each sample path corresponds to one possible outcome of the random
experiment.
\end{frame}

\begin{frame}{Example of a Sample Path}
\protect\phantomsection\label{example-of-a-sample-path}
Suppose

\[
X_t=
\begin{cases}
1, & \text{alive},\\
0, & \text{dead}.
\end{cases}
\]

One possible sample path is

{\def\LTcaptype{none} % do not increment counter
\begin{longtable}[]{@{}lccccc@{}}
\toprule\noalign{}
Time \(t\) & 0 & 0.5 & 1.0 & 1.5 & 2.0 \\
\midrule\noalign{}
\endhead
\(X_t\) & 1 & 1 & 1 & 0 & 0 \\
\bottomrule\noalign{}
\end{longtable}
}

This individual survives until sometime between \(t=1.0\) and \(t=1.5\),
after which the process remains in state 0.
\end{frame}

\begin{frame}{Summary of the Components}
\protect\phantomsection\label{summary-of-the-components}
Every stochastic process can be described by four components.

{\def\LTcaptype{none} % do not increment counter
\begin{longtable}[]{@{}ll@{}}
\toprule\noalign{}
Component & Description \\
\midrule\noalign{}
\endhead
Probability space & Defines the underlying random experiment. \\
Index set & Specifies when the process is observed. \\
State space & Specifies the possible values of the process. \\
Sample path & Describes one possible evolution of the process. \\
\bottomrule\noalign{}
\end{longtable}
}

These components provide a common framework for studying different types
of stochastic processes.
\end{frame}

\begin{frame}{Classification of Stochastic Processes}
\protect\phantomsection\label{classification-of-stochastic-processes}
Stochastic processes can be classified according to

\begin{itemize}
\tightlist
\item
  \textbf{Time domain} --- whether time is discrete or continuous.
\item
  \textbf{State space} --- whether the possible values are discrete or
  continuous.
\end{itemize}

Combining these two characteristics gives four broad classes of
stochastic processes.

These classes provide a useful framework for selecting an appropriate
model.
\end{frame}

\begin{frame}{Classification by Time Domain and State Space}
\protect\phantomsection\label{classification-by-time-domain-and-state-space}
{\def\LTcaptype{none} % do not increment counter
\begin{longtable}[]{@{}
  >{\raggedright\arraybackslash}p{(\linewidth - 4\tabcolsep) * \real{0.2203}}
  >{\raggedright\arraybackslash}p{(\linewidth - 4\tabcolsep) * \real{0.3729}}
  >{\raggedright\arraybackslash}p{(\linewidth - 4\tabcolsep) * \real{0.4068}}@{}}
\toprule\noalign{}
\begin{minipage}[b]{\linewidth}\raggedright
Time Domain
\end{minipage} & \begin{minipage}[b]{\linewidth}\raggedright
Discrete State Space
\end{minipage} & \begin{minipage}[b]{\linewidth}\raggedright
Continuous State Space
\end{minipage} \\
\midrule\noalign{}
\endhead
\textbf{Discrete time} & Discrete-time Markov chainsNCD system & Random
walksAutoregressive models \\
\textbf{Continuous time} & Poisson processesBirth-death
processesContinuous-time Markov chains & Brownian motionDiffusion
processes \\
\bottomrule\noalign{}
\end{longtable}
}

This classification provides a roadmap for the topics studied in this
course.
\end{frame}

\begin{frame}{Discrete-Time vs Continuous-Time}
\protect\phantomsection\label{discrete-time-vs-continuous-time}
The \textbf{time domain} specifies when the process is observed.

\begin{itemize}
\item
  \textbf{Discrete-time}

  Observations are made at distinct time points.

  Example: NCD levels recorded once every policy year.
\item
  \textbf{Continuous-time}

  The process is defined for every instant.

  Examples: Claim arrivals and survival status.
\end{itemize}

The choice depends on how the system evolves over time.
\end{frame}

\begin{frame}{Discrete-State vs Continuous-State}
\protect\phantomsection\label{discrete-state-vs-continuous-state}
The \textbf{state space} specifies the possible values of the process.

\begin{itemize}
\item
  \textbf{Discrete-state}

  The process takes countable values.

  Examples: NCD level, claim count, survival status.
\item
  \textbf{Continuous-state}

  The process can take any value within an interval.

  Examples: Stock prices and interest rates.
\end{itemize}

Different applications require different state spaces.
\end{frame}

\begin{frame}{Choosing an Appropriate Model}
\protect\phantomsection\label{choosing-an-appropriate-model}
The choice of stochastic process depends on the characteristics of the
system.

Ask the following questions.

\begin{enumerate}
\item
  \textbf{When is the process observed?}

  Is time discrete or continuous?
\item
  \textbf{What values can the process take?}

  Is the state space discrete or continuous?
\item
  \textbf{How does the process evolve?}

  Which stochastic model best describes its behaviour?
\end{enumerate}

These questions guide the selection of an appropriate stochastic
process.
\end{frame}

\begin{frame}{Roadmap of This Course}
\protect\phantomsection\label{roadmap-of-this-course}
The remainder of the course studies several important classes of
stochastic processes.

{\def\LTcaptype{none} % do not increment counter
\begin{longtable}[]{@{}ll@{}}
\toprule\noalign{}
Topic & Process Type \\
\midrule\noalign{}
\endhead
Discrete-time Markov chains & Discrete time, discrete state \\
Continuous-time Markov processes & Continuous time, discrete state \\
Poisson processes & Continuous time, discrete state \\
Brownian motion & Continuous time, continuous state \\
\bottomrule\noalign{}
\end{longtable}
}

Each topic develops techniques for analysing the behaviour of stochastic
processes.
\end{frame}

\begin{frame}{Chapter Summary}
\protect\phantomsection\label{chapter-summary}
In this chapter, we introduced the basic ideas of stochastic processes.

You should now be able to

\begin{itemize}
\tightlist
\item
  \textbf{Explain} why stochastic processes are used to model systems
  that evolve over time.
\item
  \textbf{Distinguish} between a random variable and a stochastic
  process.
\item
  \textbf{Identify} the stochastic process, time domain, and state space
  in simple actuarial applications.
\item
  \textbf{Interpret} the notation \[
  \{X_t:t\in T\}.
  \]
\item
  \textbf{Describe} the main components of a stochastic process.
\item
  \textbf{Classify} stochastic processes according to their time domain
  and state space.
\end{itemize}

These ideas form the foundation for the remainder of the course.
\end{frame}




\end{document}
